\begin{abstract}
Ο υπερτουρισμός αποτελεί ένα από τα πιο πιεστικά ζητήματα των σύγχρονων
κοινωνιών, καθώς η υπέρβαση της φέρουσας ικανότητας ενός προορισμού
επηρεάζει αρνητικά τόσο τους μόνιμους κατοίκους όσο και την ίδια την
ποιότητα της τουριστικής εμπειρίας. Τα παραδοσιακά στατιστικά εργαλεία,
όπως οι δείκτες τουριστικής έντασης και πυκνότητας, αδυνατούν συχνά να
αποτυπώσουν τη δυναμική που προκύπτει από την αλληλεπίδραση κατοίκων και
τουριστών μέσα στον αστικό χώρο.

Η παρούσα διπλωματική εργασία αναπτύσσει ένα μοντέλο βασισμένο σε πράκτορες
\en{(Agent-Based Model)} για την προσομοίωση του φαινομένου του υπερτουρισμού
σε πραγματικό αστικό περιβάλλον, με περιοχή μελέτης την πόλη της Πάτρας.
Το οδικό δίκτυο της πόλης κατασκευάζεται ως γράφος με αξιοποίηση ελεύθερα
προσβάσιμων γεωχωρικών δεδομένων του \en{OpenStreetMap}, ενώ πάνω σε αυτό
μοντελοποιούνται δύο τύποι πρακτόρων: οι μόνιμοι κάτοικοι, οι οποίοι
παραμένουν στάσιμοι και λειτουργούν ως δείκτες τουριστικής πίεσης μέσω της
μεταβλητής ευτυχίας τους, και οι τουρίστες, οι οποίοι κινούνται έξυπνα στο
οδικό δίκτυο ακολουθώντας ιεραρχημένη λογική προτεραιοτήτων τριών επιπέδων
(αξιοθέατα, εστιατόρια/καφέ, καταστήματα), εμπλουτισμένη με μηχανισμό
εμπρόσθιας οπτικής σάρωσης για ευκαιριακές παρεκτροπές.

Η υλοποίηση πραγματοποιείται σε περιβάλλον ανοιχτού κώδικα με τη γλώσσα
\en{Python} και το \en{framework Mesa}, ενώ η οπτικοποίηση των αποτελεσμάτων
γίνεται μέσω διαδραστικού \en{web} χάρτη βασισμένου στη βιβλιοθήκη
\en{Leaflet}, ο οποίος επιτρέπει την παρακολούθηση της προσομοίωσης σε
πραγματικό χρόνο, καθώς και τη χαρτογράφηση θερμικών χαρτών θορύβου και
ρύπανσης και γραφημάτων κατανομής ευτυχίας και ικανοποίησης.

Η αξιολόγηση του μοντέλου πραγματοποιήθηκε μέσω τριών σεναρίων προσομοίωσης
διαφορετικής τουριστικής πίεσης. Τα αποτελέσματα αναδεικνύουν ότι η επίδραση
του τουρισμού δεν κατανέμεται ομοιόμορφα στον χώρο, αλλά συγκεντρώνεται γύρω
από συγκεκριμένες ζώνες υψηλής τουριστικής κυκλοφορίας, ενώ η πλειονότητα των
κατοίκων και των τουριστών διατηρεί ικανοποιητικά επίπεδα ευημερίας ακόμα και
σε σενάρια αυξημένης τουριστικής πίεσης. Η γεωγραφική γενικευσιμότητα του
συστήματος του επιτρέπει να εφαρμοστεί άμεσα σε οποιαδήποτε ελληνική πόλη ή
χωριό, προσφέροντας ένα αναπαραγώγιμο εργαλείο ανάλυσης και διαχείρισης του
φαινομένου του υπερτουρισμού.
\end{abstract}

\begin{abstracteng}
\en{Overtourism has become one of the most pressing issues in modern societies,
as exceeding a destination's carrying capacity negatively affects both permanent
residents and the quality of the tourist experience itself. Traditional
statistical tools, such as tourist intensity and density indices, often fail
to capture the dynamics that emerge from the interaction between residents and
tourists within an urban environment.}

\en{This diploma thesis develops an Agent-Based Model (ABM) for simulating the
phenomenon of overtourism in a real urban environment, using the city of Patras
as the study area. The city's road network is constructed as a graph using freely
available geospatial data from OpenStreetMap, on top of which two types of agents
are modeled: permanent residents, who remain stationary and act as indicators of
tourist pressure through their happiness variable, and tourists, who move
intelligently across the road network following a three-tier hierarchical priority
logic (attractions, restaurants/cafes, shops), enriched with a forward vision scan
mechanism for opportunistic detours.}

\en{The implementation is carried out in an open-source environment using the
Python programming language and the Mesa framework, while the results are
visualized through an interactive web map based on the Leaflet library, allowing
real-time monitoring of the simulation as well as the mapping of noise and
pollution heatmaps and charts of happiness and satisfaction distribution.}

\en{The model was evaluated through three simulation scenarios with different levels
of tourist pressure. The results show that the impact of tourism is not uniformly
distributed across space, but rather concentrated around specific zones of high
tourist traffic, while the majority of residents and tourists maintain
satisfactory wellbeing levels even under scenarios of increased tourist pressure.
The geographic generalizability of the system allows it to be directly applied to
any Greek city or village, offering a reproducible tool for analyzing and managing
the phenomenon of overtourism.}
\end{abstracteng}